% ==============================================================================
% Fichier : front/03b_labs.tex : SAFI (SIS589)
% ==============================================================================

\section*{Les dix labs de NDAP SUKULU}

\begin{tableaucours}{L{2.4cm}L{8.4cm}L{1.1cm}}{%
    \tableauHeader{Lab} &
    \tableauHeader{Sujet} &
    \tableauHeader{Ch.}
}
SAFI-LAB-01 & Lire un tableau de bord de supervision & 3 \\
SAFI-LAB-02 & Analyser des journaux et repérer les angles morts & 4 \\
SAFI-LAB-03 & Écrire et éprouver une règle de détection (SIEM) & 7 \\
SAFI-LAB-04 & Reconnaître une intrusion dans des alertes corrélées & 5 \\
SAFI-LAB-05 & Chaîne de custody et qualification d'artefacts & 11 \\
SAFI-LAB-06 & Forensique système~: reconstituer une chronologie & 12 \\
SAFI-LAB-07 & Forensique réseau~: exfiltration et pilotage (flux exportés) & 14 \\
SAFI-LAB-08 & Classer et traiter des incidents (gravité, délais) & 15 \\
SAFI-LAB-09 & Mapping d'un incident à MITRE ATT\&CK & 17 \\
SAFI-LAB-10 & Couverture de détection et verdict OUI/NON/NEK & 20 \\
\end{tableaucours}

\medskip
Chaque lab dure environ une heure et se corrige automatiquement. Les pièces
se trouvent dans le dossier \texttt{preuves} (lecture seule) et une copie de
travail dans \texttt{travail}.
